\chapter*{Conclusions of Part~III}
\addcontentsline{toc}{chapter}{Conclusions of Part~III}
\label{ch:part3-concl}

The three chapters in this Part establish a comprehensive framework for extracting the neutrino-induced phase shift from cosmological data, progressing from its high-significance detection in the CMB to direct constraints on neutrino interactions and, finally, to the identification of a novel signature in the 21-cm spectrum from cosmic dawn.

Specifically, in Chapter~\ref{ch:neutrino-phase-shift}, I introduced two complementary template-based methods for isolating the phase shift in CMB data. The spectrum-based template measures the coherent displacement of acoustic peaks across multipoles, while the perturbation-based template captures the shift directly at the level of photon-baryon perturbations where it is originally imprinted. Applying both methods to current CMB datasets, I detected a non-zero neutrino-induced phase shift with a $\sim 14\sigma$ significance and confirmed consistency with the Standard Model expectation of three free-streaming neutrino species. The perturbation-based method proved particularly robust across different data releases, and forecasts for the Simons Observatory indicate that the precision on the phase shift will improve by roughly 50\%  relative to current measurements, reaching a sensitivity at which deviations from the Standard Model prediction could begin to be meaningfully probed.

Chapter~\ref{ch:neutrino-interactions} exploited this framework to place direct constraints on neutrino interactions. By modeling two physically motivated classes of interaction scenarios with scattering rates scaling as $T_\nu^3$ and $T_\nu^5$, where $T_\nu$ is the neutrino temperature, I showed that the phase-shift amplitude maps cleanly onto the neutrino decoupling redshift, establishing lower bounds of $z_{\nu,\mathrm{dec}} > 1.3 \times 10^4$ and $z_{\nu,\mathrm{dec}} > 1.7 \times 10^4$ respectively from combined CMB data. These constraints firmly require that neutrinos were free-streaming well before matter-radiation equality and are competitive with those obtained from dedicated model-dependent analyses. This highlights the unique power of the phase shift as a clean theoretical lever for testing neutrino properties, and more broadly showcases the effectiveness of signature-driven approaches in extracting fundamental physics from cosmological data.

Finally, in Chapter~\ref{ch:neutrino-21cm}, I extended the phase-shift analysis to the 21-cm brightness temperature during cosmic dawn, identifying a qualitatively new feature. The 21-cm signal contains both BAOs and VAOs, and the neutrino-induced phase shift in each component differs in both amplitude and scale dependence. The VAO phase shift asymptotes to a value roughly 1.5 times larger than the BAO phase shift and exhibits a smoother transition between large and small scales, reflecting its origin in four-point rather than two-point correlation functions. The combined signature evolves across cosmic dawn as the relative contributions of VAOs and BAOs shift with astrophysical conditions, producing a distinctive redshift-dependent signal that is robust at small scales but sensitive to star formation physics on large scales. By fully characterizing this signal, including its dependence on astrophysical parameters, this work lays the foundation for extracting the neutrino-induced phase shift from upcoming 21-cm observations with HERA~\citep{deboer17,hera23} and SKA~\citep{SKA:2018ckk, SKA19}, opening a new and independent window into the free-streaming nature of neutrinos.

Taken together, the work presented in these three chapters advances the neutrino-induced phase shift from a clean confirmation of the cosmic neutrino background in the CMB to a precision probe of neutrino interactions and early-Universe physics across multiple cosmological observables. As current and upcoming surveys deliver increasingly precise measurements of both CMB anisotropies and large-scale structure probes --- including the 21-cm signal from cosmic dawn --- this robust signature will continue to provide unique insights into the fundamental properties of neutrinos and their potential interactions with the particle content of the early Universe.
